\section{Conditional Trajectory Abstraction}
\label{sec:method}

\method has three learned components: a target encoder that turns the actual future into a code (\cref{sec:source}), two heads that read the code (\cref{sec:reader}), and a world model that predicts the code from an action chunk (\cref{sec:wm}).
We describe them, the training procedure (\cref{sec:codesign}), and deployment (\cref{sec:planning}) in turn.

\subsection{Inputs}
\label{sec:context}
As in DINO-WM~\cite{zhou2025dinowm}, the image encoder $f$ is a frozen DINOv2 ViT-S/14~\cite{oquab2024dinov2}; we project its patch features to 128 dimensions with PCA fitted on training frames.
The context $C$ contains the $16{\times}16$ patch tokens of the current image, the previous image pooled to $8{\times}8$ tokens, and the proprioception at both times.
The future segment $\tau$ contains the full token grid of the final image, the pooled tokens of three intermediate images with temporal embeddings, and the final proprioception.
Features are computed once per image and shared by all candidates and goals.

\subsection{Target encoder}
\label{sec:source}
The target encoder $q_\phi(C,\tau)$ is used only during training.
$M$ learned queries cross-attend to the tokens of the future and of the context, and each output is quantized with FSQ~\cite{mentzer2024fsq} using levels $(8,8,4)$, so that each token takes one of $V{=}256$ values.
Because the queries also attend to the context, the encoder can learn which parts of the future the history already implies and spend the code on the rest.
The $M$ tokens describe the segment as a whole rather than one frame each.

\subsection{Reader and feature decoder}
\label{sec:reader}
The reader $D_\psi(C,S,g)$ is a Transformer over the context tokens, the code tokens, and the patch tokens of a goal image, and it outputs a single score.
It is trained on sibling candidates that branch from the same state, with a pairwise logistic loss over the sibling pairs whose labels differ by more than a margin $\epsilon$:
\begin{equation}
  \mathcal{L}_{\mathrm{rank}}=\E_{(i,j):\,\ell_i>\ell_j+\epsilon}\Big[w_{ij}\log\big(1+e^{-(s_i-s_j)}\big)\Big],
\label{eq:rank}
\end{equation}
where $s_i=D_\psi(C,S_i,g)$ and $w_{ij}{=}1$ when training the reader.
When the loss is applied to predicted codes (\cref{sec:wm}), $w_{ij}=\min\big(1,(\ell_i-\ell_j)/\sigma\big)$ down-weights pairs whose labels barely differ.
Ranking within a decision matches how the score is used, so the score does not need to be calibrated across states.

A ranking loss alone would let the code shrink to whatever orders the candidates for the training goals.
The feature decoder $G_\xi(C,S)$ prevents this by reconstructing the future tokens from the context and the code.
It predicts the change of each token relative to the current image, and its squared-error loss $\mathcal{L}_{\mathrm{feat}}$ is normalized by the error of predicting no change, so that a code that carries no information scores one.
The decoder thus plays the role that EMA targets and variance regularizers play in other JEPAs~\cite{assran2023ijepa,lecun2022path}: it keeps the target from degenerating, here into a score for the training goals.
No image is ever generated.

\subsection{World model}
\label{sec:wm}
The world model $p_\theta(S\mid C,\bA)$ is the predictor of the JEPA.
A Transformer encoder reads the context tokens and one token per action, and $M$ learned queries decode all code tokens in parallel.
Each FSQ coordinate of each token is predicted as a categorical distribution over its levels, so one forward pass yields a distribution over codes for every candidate in the batch.
The reader receives the expected codeword $\hat S=\E_{p_\theta}[S]$, which exposes the world model's uncertainty to the reader and avoids the ties that the most likely code produces among similar candidates (\cref{sec:ablations}).
The 128-bit bound therefore applies to the target code; the deployed input to the reader is the predictive mean, $M{\times}3$ numbers in the FSQ cube.
The world model never sees future images, the goal, or labels.

The target encoder and the reader are frozen while the world model is trained against the stop-gradient code $S=\sg\big(q_\phi(C,\tau)\big)$:
\begin{equation}
  \mathcal{L}_{\mathrm{WM}}=\mathcal{L}_{\mathrm{NLL}}+\mathcal{L}_{\mathrm{cons}}+\mathcal{L}^{\hat S}_{\mathrm{rank}}.
  \label{eq:wm}
\end{equation}
$\mathcal{L}_{\mathrm{NLL}}$ is the cross-entropy of each coordinate of $S$.
The other two terms pass gradients through the frozen reader into the world model.
$\mathcal{L}_{\mathrm{cons}}$ matches the within-decision differences between the reader's scores of $\hat S$ and of $S$, and $\mathcal{L}^{\hat S}_{\mathrm{rank}}$ applies the weighted form of \eqref{eq:rank} to the scores of the predicted codes.
A context-only model $r_\omega(S\mid C)$ is trained with $\mathcal{L}_{\mathrm{NLL}}$ alongside the world model.
Its excess cross-entropy over the world model estimates $I(S;\bA\mid C)$, the part of the code that the action explains.

\subsection{Training}
\label{sec:codesign}

\paragraph{Branched data.}
We roll out the frozen policy and, at every decision, execute each of its $K$ proposals for $\He$ steps from an exact copy of the simulator state.
This yields $K$ sibling futures and labels that share one history.
The proposals of a competent policy often lead to nearly identical futures, so we also branch a perturbed copy of each proposal, shifted by a constant planar offset whose scale grows across the copies.
To cover the states that a selector visits, not only those that the policy visits, some rollouts execute the best-labeled proposal at a random half of their decisions.
All branches of a rollout stay in the same data split.

\paragraph{Stages.}
Stage~1 trains the target encoder, the reader, and the feature decoder on actual futures with $\mathcal{L}_{\mathrm{code}}=\mathcal{L}_{\mathrm{rank}}+\alpha\,\mathcal{L}_{\mathrm{feat}}$, using straight-through gradients through the quantizer.
Stage~2 freezes them and trains the world model with \eqref{eq:wm}.
Each network keeps the checkpoint with the best value of its own objective on held-out development rollouts.
Throughout training we monitor the code perplexity, the number of distinct codes per decision, the decoder's $R^2$, and $I(S;\bA\mid C)$ to detect collapse.

\paragraph{Co-design (ablation).}
The code can also be shaped to be predictable.
An optional third stage alternates blocks of target-encoder and world-model updates and adds $\lambda\,\E\big[-\log \bar p_\theta(S\mid C,\bA)\big]$ to $\mathcal{L}_{\mathrm{code}}$, where $\bar p_\theta$ is a frozen copy of the world model and the discrete likelihood is interpolated multilinearly on the FSQ grid to obtain a gradient.
We report it as an ablation (\cref{sec:ablations}).

\subsection{Planning}
\label{sec:planning}
\Cref{alg:cta} summarizes deployment.
At each decision, \method encodes the images once, runs the world model once on all proposals, and runs the reader on every candidate--goal pair.
A new goal therefore costs only reader evaluations.

\begin{algorithm}[t]
\caption{Selecting among policy proposals with \method.}
\label{alg:cta}
\small
\begin{algorithmic}[1]
\Require frozen policy $\pi$ and image encoder $f$; world model $p_\theta$; reader $D_\psi$; goal images $\cG$
\While{episode not finished}
  \State $C\gets$ context from $f$ applied to the history \Comment{once per decision}
  \State $\{\bA^{(k)}\}_{k=1}^{K}\gets$ proposals of $\pi$
  \State $\hat S^{(k)}\gets\E_{p_\theta(\cdot\mid C,\bA^{(k)})}[S]$ for all $k$ \Comment{one batched pass}
  \State $s_k\gets\frac{1}{|\cG|}\sum_{g\in\cG}D_\psi\big(C,\hat S^{(k)},g\big)$
  \State $k^\star\gets$ smallest index attaining $\max_k s_k$ \Comment{ties keep $\bA^{(1)}$}
  \State execute $\bA^{(k^\star)}$; $t\gets t+\He$
\EndWhile
\end{algorithmic}
\end{algorithm}

\subsection{Implementation details}
\label{sec:impl}
All modules have width 256 and 8 attention heads.
The target encoder has three Transformer decoder layers, the reader four encoder layers, the feature decoder two decoder layers, and the world model four encoder and four decoder layers.
We use $M{=}16$ (128 bits per target code), $\alpha{=}1$, $\epsilon{=}10^{-3}$, and $\sigma{=}10^{-2}$.
A small penalty on FSQ pre-activations beyond $\pm1.5$ keeps the quantizer out of saturation, where straight-through gradients vanish.
Stage~2 runs 8{,}000 AdamW updates with learning rate $10^{-4}$, weight decay $0.05$, and bfloat16 autocast, with 32 decisions per update.
\TODO{Final stage-1 schedule, trained from scratch; current checkpoints are warm-started across development rounds.}
On PushT, the $96{\times}96$ images are resized to $224{\times}224$ for DINOv2, $\He{=}8$, and the intermediate images follow steps 2, 4, and 6 of the chunk.
Each action is a target position of the agent and is embedded in absolute coordinates and relative to the agent's current position.
